\subsubsection{Analisis Gap Infrastruktur Audit Trail terhadap DecMed}

Setelah komponen-komponen utama infrastruktur \textit{audit trail} diidentifikasi berdasarkan model konseptual ISO 27789, langkah selanjutnya adalah membandingkan model tersebut dengan kondisi eksisting sistem DecMed. Analisis ini bertujuan untuk mengidentifikasi komponen infrastruktur \textit{audit trail} yang telah tersedia maupun komponen yang masih belum tersedia pada arsitektur DecMed. Hasil analisis tersebut digunakan untuk menentukan kebutuhan implementasi yang harus dipenuhi agar DecMed mampu mendukung proses pencatatan, penyimpanan, pemantauan, dan penyajian \textit{audit trail} secara komprehensif.

Analisis dilakukan dengan mengevaluasi setiap komponen yang direkomendasikan oleh ISO 27789 terhadap arsitektur DecMed yang terdiri atas aplikasi klien, \textit{smart contract} IOTA Move, \textit{Proxy Re-Encryption Server}, serta infrastruktur penyimpanan data yang telah ada. Hasil analisis ditunjukkan pada Tabel \ref{tab:gap_komponen_audit}.

\begin{xltabular}{\textwidth}{|p{3.2cm}|p{3.2cm}|X|X|}
\caption{Analisis Gap Infrastruktur Audit Trail terhadap DecMed}
\label{tab:gap_komponen_audit}\\

\hline
\textbf{Komponen} &
\textbf{Kondisi DecMed} &
\textbf{Analisis Gap} &
\textbf{Kebutuhan}\\
\hline
\endfirsthead

\hline
\textbf{Komponen} &
\textbf{Kondisi DecMed} &
\textbf{Analisis Gap} &
\textbf{Kebutuhan}\\
\hline
\endhead

\hline
\endfoot

Audit Logger Service &
Belum tersedia. Audit event dicatat secara terpisah oleh masing-masing komponen. &
Belum terdapat mekanisme terpusat untuk menerima seluruh audit event dari berbagai komponen sistem. &
Diperlukan layanan yang bertugas menerima seluruh audit event sebelum diproses lebih lanjut.\\
\hline

Audit Record Generator Service &
Belum tersedia. &
Belum terdapat mekanisme pembentukan audit record yang konsisten sehingga struktur audit record bergantung pada implementasi masing-masing komponen. &
Diperlukan mekanisme pembentukan audit record yang memiliki format dan atribut yang seragam.\\
\hline

Audit Event Catalogue Service &
Belum tersedia. &
Belum terdapat definisi terpusat mengenai jenis audit event beserta atribut yang dimiliki. &
Diperlukan katalog audit event sebagai acuan seluruh komponen dalam menghasilkan audit record.\\
\hline

Audit Repository &
Belum tersedia. &
DecMed belum memiliki repositori khusus untuk menyimpan audit trail secara terpusat. Riwayat akses pasien dan transaksi blockchain belum dapat berfungsi sebagai repositori audit trail yang komprehensif. &
Diperlukan repositori khusus yang mampu menyimpan audit trail secara terpusat dan mendukung kebutuhan investigasi.\\
\hline

Audit Monitor Service &
Belum tersedia. &
Belum terdapat mekanisme pemantauan terhadap pola aktivitas atau kejadian mencurigakan secara otomatis. &
Diperlukan mekanisme monitoring audit event sebagai dasar deteksi aktivitas tertentu.\\
\hline

Audit Alert Service &
Belum tersedia. &
Belum terdapat mekanisme pemberian notifikasi ketika ditemukan aktivitas yang memenuhi kondisi tertentu. &
Diperlukan layanan yang mampu menghasilkan notifikasi berdasarkan hasil pemantauan audit trail.\\
\hline

Audit Report Service &
Sebagian tersedia melalui \textit{Patient Access Log}. &
Riwayat akses yang tersedia hanya mencakup aktivitas tertentu dan belum mendukung pencarian maupun pelaporan audit trail secara menyeluruh. &
Diperlukan layanan pelaporan yang mampu melakukan pencarian dan penyajian audit trail untuk proses audit maupun investigasi.\\
\hline

\end{xltabular}

Berdasarkan Tabel \ref{tab:gap_komponen_audit}, dapat dilihat bahwa sebagian besar komponen infrastruktur \textit{audit trail} yang direkomendasikan oleh ISO 27789 belum tersedia pada arsitektur DecMed. Meskipun DecMed telah memiliki beberapa mekanisme pencatatan aktivitas, seperti \textit{Patient Access Log} dan riwayat transaksi pada IOTA DLT, mekanisme tersebut belum membentuk sebuah sistem \textit{audit trail} yang terintegrasi. Selain itu, belum tersedia layanan yang secara khusus bertanggung jawab untuk menerima \textit{audit event}, membentuk \textit{audit record}, menyimpan data audit secara terpusat, maupun menyediakan fasilitas pemantauan dan pelaporan.

Hasil analisis ini menunjukkan bahwa pengembangan sistem \textit{audit trail} pada DecMed tidak hanya memerlukan penambahan media penyimpanan audit trail, tetapi juga memerlukan beberapa komponen layanan baru yang bekerja secara terintegrasi. Komponen-komponen tersebut selanjutnya akan dianalisis berdasarkan kontrol keamanan yang direkomendasikan oleh NIST SP 800-53 Revision 5 untuk menentukan kebutuhan fungsional dan non-fungsional yang harus dipenuhi oleh setiap komponen sebelum dilakukan proses perancangan sistem.